\begin{table}[htbp]
\centering
\caption{Standardized Effect Sizes}
\label{tab:sde}
\begin{adjustbox}{max width=\textwidth}
\begin{tabular}{lcccccc}
\hline\hline
Outcome & $\hat{\beta}$ & SE & SD($Y$) & SDE & SE(SDE) & Classification \\
\hline
\multicolumn{7}{l}{\textit{Panel A: Pooled}} \\
Bunching ratio (all bridges) & 0.077 & 0.010 & 0.027 & 2.880 & 0.387 & Large positive \\
[6pt]
\multicolumn{7}{l}{\textit{Panel B: Heterogeneous (by owner type)}} \\
State DOT bridges & 0.004 & --- & 0.041 & 0.091 & --- & Moderate positive \\
Local government bridges & 0.018 & --- & 0.040 & 0.448 & --- & Large positive \\
\hline\hline
\end{tabular}
\end{adjustbox}
\begin{tablenotes}[flushleft]\footnotesize
\item \textit{Notes:}
\textbf{Country:} United States.
\textbf{Research question:} Does the elimination of sufficiency-based federal bridge funding (MAP-21, 2012) reduce strategic manipulation of bridge condition ratings by state transportation agencies?
\textbf{Policy mechanism:} The Highway Bridge Program apportioned federal replacement funds to states based on the count of bridges with sufficiency ratings below 50; MAP-21 replaced this formula with a performance-based system using structurally deficient deck-area triggers, weakening the incentive to deflate ratings below the threshold.
\textbf{Outcome definition:} Bunching ratio, defined as the count of bridges with integer sufficiency ratings in [46,49] divided by [50,53], measuring excess mass just below the federal funding threshold.
\textbf{Treatment:} Binary; pre-MAP-21 (2000--2012) versus post-MAP-21 (2013--2018).
\textbf{Data:} FHWA National Bridge Inventory, 2000--2018, bridge-year panel, approximately 11.5 million observations across 50 states and DC.
\textbf{Method:} Polynomial bunching estimator (order 7) with difference-in-bunching pre/post MAP-21; bootstrap standard errors (200 replications, resampled by bridge ID).
\textbf{Sample:} All public bridges in the 50 US states and DC with non-missing sufficiency ratings; excludes territories.
SDE $= \hat{\beta} / \text{SD}(Y)$ where SD($Y$) is the pre-treatment standard deviation of the annual bunching ratio. Classification refers to magnitude, not statistical significance: Large ($|$SDE$| > 0.15$), Moderate ($0.05$--$0.15$), Small ($0.005$--$0.05$), Null ($< 0.005$).
\end{tablenotes}
\end{table}

